% =====================================================================
%  Conferenze, Zagarolo (On The Road), giorno 2 — Intervista 1
%  L'importanza degli eventi internazionali per gli studi indipendenti
%  Fonte: trascrizione automatica "Track 15". Intervista in inglese a
%  Diego, business development manager di Games London (London Games
%  Festival), condotta dal curatore dell'area indie di un festival
%  italiano, con domande dal pubblico.
%  Etichette con prefisso "c3-1:". Nessun pacchetto aggiuntivo.
% =====================================================================

\chapter[Eventi internazionali per gli studi indie]{L'importanza degli eventi
internazionali per gli studi indipendenti}
\label{cap:c3-1-eventi-internazionali}

Il secondo giorno di conferenze si apre con un'intervista dedicata agli
\textbf{eventi} del settore: fiere per il pubblico, conferenze professionali,
mercati B2B. Quali scegliere, come presentarsi, che cosa fare durante un incontro
con un publisher e dopo l'evento, quando invece è meglio restare a casa a lavorare
sul gioco. Il capitolo è organizzato così:
\begin{enumerate}
  \item l'ospite e Games London (sezione~\ref{sec:c3-1-ospite});
  \item come scegliere gli eventi (sezione~\ref{sec:c3-1-scegliere});
  \item come si viene selezionati: il London Games Festival
        (sezione~\ref{sec:c3-1-selezione});
  \item trasformare un incontro in un'opportunità
        (sezione~\ref{sec:c3-1-incontro});
  \item dopo l'evento: il follow-up (sezione~\ref{sec:c3-1-dopo});
  \item quando non andare a un evento (sezione~\ref{sec:c3-1-quando-no});
  \item le domande del pubblico (sezione~\ref{sec:c3-1-domande}).
\end{enumerate}

% ---------------------------------------------------------------------
\section{L'ospite}
\label{sec:c3-1-ospite}

Il moderatore presenta l'ospite come ``collega e amico'', \textbf{business
development manager} di \textbf{Games London}, l'organizzazione che sta dietro al
\textbf{London Games Festival}\footnote{Nella trascrizione l'ospite è presentato
come ``Livio Grammatico'', ma subito dopo il moderatore gli dà la parola chiamandolo
``Diego'', nome che ricorre poi per tutta l'intervista e con cui lo ringrazia in
chiusura. Adottiamo ``Diego''; non è chiaro se il nome iniziale si riferisca a
un'altra persona o sia un errore di trascrizione. L'organizzazione è trascritta
anche come ``GameStone'', il festival come ``Ronald Rex Festival'' e ``Rondon
Essential''.}. Il tema è \textbf{l'importanza degli eventi internazionali} per gli
studi di sviluppo: perché partecipare e che cosa si può ottenere parlando con gli
operatori internazionali del settore.

Diego descrive Games London come un'\textbf{organizzazione culturale} sostenuta dal
\textbf{sindaco di Londra}\footnote{Nella trascrizione ``the mayor of Loma'',
corretto subito dopo dallo stesso relatore.} e, da quest'anno, dal \textbf{governo
britannico}, con lo scopo di far crescere l'ecosistema dei videogiochi nel Regno
Unito e nella capitale. Organizza il festival e programmi di supporto agli studi,
britannici e stranieri, per dare loro le competenze e gli strumenti per crescere e
\textbf{accedere ai mercati}.

% ---------------------------------------------------------------------
\section{Come scegliere gli eventi}
\label{sec:c3-1-scegliere}

\subsection{Il ritorno sull'investimento}

Gli eventi nel mondo sono moltissimi: showcase, fiere, conferenze. Ma partecipare è
sempre una \textbf{spesa} di viaggio, tempo, risorse ed energia; per questo bisogna
ragionare sul \textbf{ritorno sull'investimento}.

\subsection{Partire dal locale}

Uno studio emergente, al primo o secondo gioco, dovrebbe partire da ciò che succede
nella \textbf{propria industria locale}: gli eventi più vicini e più comodi da
raggiungere. Diego cita una conversazione del giorno prima con
Dario\footnote{La trascrizione riporta ``Dario from AAP there''; con ogni
probabilità si tratta di IIDEA, l'associazione dell'industria italiana, che a
Firenze organizza l'evento professionale \emph{First Playable}. La ricostruzione è
un'ipotesi.} su un evento a \textbf{Firenze}: per uno sviluppatore italiano
all'inizio del percorso sarebbe il primo evento da considerare.

\subsection{Fiere per il pubblico e conferenze professionali}

Bisogna distinguere due tipi di eventi, importanti per ragioni molto diverse.
\begin{description}
  \item[Fiere per il pubblico (\emph{consumer show}).] L'obiettivo è far vedere e
        provare il gioco al \textbf{maggior numero di persone} nel minor tempo
        possibile. È fondamentale che l'evento abbia anche una
        \textbf{componente online}, perché oggi gran parte della visibilità
        arriva da lì: molti eventi fisici hanno anche uno \textbf{showcase su
        Steam}, ed è spesso da lì che arriva l'attenzione. Meglio dare priorità
        agli eventi con una componente online.
  \item[Conferenze professionali e B2B.] Le domande da porsi sono due:
        \begin{itemize}
          \item tratta argomenti che non conosco, da cui posso \textbf{imparare}
                qualcosa di nuovo?
          \item c'è un modo, possibilmente \textbf{strutturato}, per entrare in
                contatto con publisher, investitori e partner che non conosco?
        \end{itemize}
        Alcuni eventi puntano sul networking \textbf{informale}: bisogna essere
        molto sicuri di sé e avvicinare qualcuno con un bicchiere di vino o un
        caffè. Altri offrono una \textbf{piattaforma di incontri} su cui inviare
        richieste in anticipo, un modo più semplice per incontrare persone, che
        Diego personalmente preferisce.
\end{description}

% ---------------------------------------------------------------------
\section{Come si viene selezionati: il London Games Festival}
\label{sec:c3-1-selezione}

\subsection{L'Official Selection}

Al London Games Festival c'è la \textbf{Official Selection}: circa \textbf{40
giochi} da tutto il mondo, in uscita, scelti perché in qualche modo
\textbf{unici} nel gameplay, nell'arte o nella creatività. Non ci sono criteri
predefiniti sul tipo di gioco: il team \textbf{prova tutto}, se c'è un prototipo
giocabile, e giudica in base alle sensazioni che il gioco trasmette.

C'è poi un evento B2B\footnote{Trascritto prima ``Blinks Minus'' e alla fine
``Vince Finance Market''; il nome corretto è con ogni probabilità \textbf{Games
Finance Market}.} che mette in contatto studi con publisher e investitori.

\subsection{Gli errori comuni nelle candidature}

Il festival riceve circa \textbf{400 candidature} l'anno e, pur non essendo un
publisher, deve valutarle rapidamente. Per questo:
\begin{itemize}
  \item deve essere chiaro \textbf{fin da subito} di che cosa parla il gioco;
  \item il \textbf{pitch deck} deve essere breve: \textbf{al massimo 10 slide}
        (arrivano deck da 30, 40, 50 slide, difficilissimi da leggere);
  \item servono \textbf{gameplay footage} e \textbf{screenshot}. Un deck fatto solo
        di parole, che racconta storia e concept senza immagini, artwork o video,
        non è un errore in senso stretto, ma rende molto difficile capire
        l'attrattiva del gioco.
\end{itemize}
Più si riesce a mostrare come il gioco \emph{appare} e come \emph{si sente}, più
rapida è la decisione. Il festival può fare domande di ritorno e cerca di essere
flessibile, ma in generale è importante saper raccontare \textbf{studio, gioco e
progetto} in modo breve e chiaro.

% ---------------------------------------------------------------------
\section{Trasformare un incontro in un'opportunità}
\label{sec:c3-1-incontro}

\subsection{Quali eventi fuori dall'Italia}

Per uno studio piccolo, Diego salterebbe gli eventi \textbf{molto grandi} come la
\textbf{Gamescom} o la \textbf{GDC}\footnote{Trascritta ``GBC''.}: c'è così tanto
rumore, e così tante aziende, che è difficilissimo emergere. Il ritorno è basso, a
meno di avere uno studio affermato, un \textbf{track record} e una rete da attivare
prima dell'evento (``sarò alla tal conferenza, ci vediamo?'').

Chi comincia fa meglio a puntare su \textbf{eventi più piccoli in Europa}, per
costruire community, rete e future collaborazioni. Oltre al London Games Festival
cita alcuni eventi ``più amichevoli''\footnote{La trascrizione è qui molto
rovinata: ``Gamescom in Berlin, which takes place in November'', ``Dead Down in
Lisbon'', ``Digital Downs in Poland''. L'ultimo è quasi certamente
\textbf{Digital Dragons} (Cracovia); gli altri due non sono ricostruibili con
sicurezza: un evento a Berlino a novembre e uno a Lisbona.}: uno a
\textbf{Berlino} a novembre, uno a \textbf{Lisbona}, e \textbf{Digital Dragons} in
Polonia.

\subsection{Durante l'incontro}

La cosa più importante, se possibile, è \textbf{far provare il gioco} durante
l'incontro: con una Steam Deck o un portatile\footnote{La trascrizione riporta
``a Steam Deck or a Wii U or a Lambo'': probabilmente ``a Steam Deck or a laptop''.},
spendere i primi \textbf{5--10 minuti} lasciando giocare l'interlocutore mentre si
spiega la logica del gioco. È tempo ben speso: si capisce subito come funziona la
meccanica e che aspetto ha il gioco.

Se non c'è una build, Diego suggerisce questa struttura:
\begin{enumerate}
  \item \textbf{gameplay footage} e immagini che diano l'idea del loop di gioco;
  \item subito dopo, \textbf{una slide di sintesi} con le caratteristiche
        essenziali e ciò che si cerca: genere, gameplay, durata prevista, data
        d'uscita, budget, motore. In questo modo si capisce subito se il publisher
        è ``sulla stessa lunghezza d'onda'', e perché;
  \item altre \textbf{5--6 slide} che approfondiscono ciascun punto.
\end{enumerate}
Alla fine: scambiarsi \textbf{email} e biglietti da visita, ringraziare e
annunciare il follow-up (``vi mando la chiave Steam e il pitch deck''), da fare
circa \textbf{una settimana dopo}.

% ---------------------------------------------------------------------
\section{Dopo l'evento: il follow-up}
\label{sec:c3-1-dopo}

Dopo circa una settimana si manda un'email cortese:
\begin{itemize}
  \item ringraziare per il tempo dedicato;
  \item allegare il \textbf{pitch deck} e una \textbf{chiave Steam}, o il link a
        un prototipo o a un trailer;
  \item chiedere un parere e proporre di restare in contatto.
\end{itemize}
Poi \textbf{aspettare}: come si è visto il giorno
prima\footnote{Si veda il capitolo~\ref{cap:c2-2-producer}.}, la pipeline di
publisher e investitori è molto lunga e possono passare settimane. Non bisogna
riscrivere dopo due giorni né ``dare la caccia'' alle persone: un \textbf{secondo
follow-up} gentile dopo due o tre settimane va bene, ma è meglio essere il più
pazienti possibile. Se la risposta non arriva, probabilmente non c'è interesse: si
passa alla persona successiva.

% ---------------------------------------------------------------------
\section{Quando non andare a un evento}
\label{sec:c3-1-quando-no}

\begin{itemize}
  \item \textbf{Fiere per il pubblico senza componente online.} Se da sviluppatore
        emergente bisogna pagare uno stand, ma l'evento non ha uno showcase o una
        promozione su Steam, conviene meno: di persona il gioco lo vedono poche
        persone, per limiti di spazio e di tempo, mentre online il pubblico
        potenziale è molto più grande. Le due cose dovrebbero andare insieme.
  \item \textbf{Eventi B2B troppo presto.} Se il pitch non è ancora curato, se non
        si sa bene che cosa chiedere o con quale budget, meglio evitare: non si
        sfrutterebbe l'occasione.
  \item \textbf{Troppi eventi.} Per uno studio indie, dopo \textbf{due, al massimo
        tre} grandi eventi l'anno il guadagno marginale di un evento in più è
        basso; anche due bastano. Il \textbf{costo opportunità} è il tempo
        sottratto al gioco. Diego vede studi che girano l'Europa per cinque, sei,
        sette, otto eventi l'anno: un'esperienza umana bellissima, che fa sentire
        meno soli, ma tempo che si poteva dedicare al progetto. Serve equilibrio.
\end{itemize}

% ---------------------------------------------------------------------
\section{Le domande del pubblico}
\label{sec:c3-1-domande}

\subsection{Vendere un meme invece di un gioco?}

Un partecipante osserva che il mercato è cambiato in modo impensabile in cinque
anni: per farsi notare sembra che gli sviluppatori debbano ``fare i pagliacci'' e
vendere un meme invece del prodotto. Che cosa è cambiato, che cosa succederà nei
prossimi cinque anni, e che ruolo avrà l'IA?

Diego non si sbilancia sulle previsioni, ma propone una riflessione. Il videogioco è
una \textbf{forma d'arte creativa} che, nella maggior parte dei casi, nasce da
\textbf{dentro di sé}, come un film o una canzone: qualcosa che si vuole condividere.
Questa visione, tratta dalla propria esperienza e dalla propria storia, rende il
gioco \textbf{unico}. Se invece nasce dalle \textbf{tendenze}, che sono visibili a
tutti, tutti seguiranno le stesse e si vedrà una moltitudine di giochi uguali; alcuni
guadagnano, ma sono la punta dell'iceberg.

Oggi molti studi non partono da qualcosa da dire, ma da domande come: come faccio un
gioco facile da \textbf{streammare su TikTok}? Che piaccia agli
\textbf{influencer}? Che un publisher \textbf{firmerà}? Devo usare l'IA? Devo fare
un gioco di un certo genere di moda anche se non mi interessa\footnote{La
trascrizione riporta ``Should I make a fence log?'': il genere citato non è
ricostruibile.}? Avere una \textbf{comprensione commerciale} è vitale (studiare il
mercato, capire quali generi vanno meglio, avere una strategia di marketing), ma si
deve sempre partire da \textbf{ciò che si ha da dire}. Seguendo le tendenze le
probabilità di essere scoperti, e di essere soddisfatti del proprio lavoro, sono
molto basse.

L'IA, per Diego, è uno \textbf{strumento} che renderà più facile fare un gioco, forse
in meno tempo e con meno persone; ma partire dall'IA quando si fa un'opera creativa
non è il punto giusto. Ripete il consiglio del giorno prima: partire da qualcosa di
\textbf{piccolo}, pubblicare il primo progetto \textbf{entro un anno}, poi passare al
successivo, senza perdere l'ispirazione che viene dalle esperienze vissute.

\subsection{Quanto deve durare una demo?}

Un partecipante chiede come bilanciare i contenuti e la durata della demo a un
evento, per far giocare più persone possibile. Diego:
\begin{itemize}
  \item ogni giocatore deve avere il tempo di sperimentare \textbf{un loop di gioco
        completo}: si parte dal loop e si costruisce una demo che permetta a
        chiunque di viverlo;
  \item meglio \textbf{non partire dall'inizio del gioco}: tra istruzioni e
        tutorial si arriva a giocare dopo un quarto d'ora. Meglio scegliere una
        \textbf{porzione più vivace}, ma con una \textbf{parte iniziale su misura}
        che spieghi comandi e obiettivi, così il giocatore non si perde;
  \item vale la pena spendere tempo su una demo \textbf{personalizzata} per
        l'evento: le persone lo apprezzano.
\end{itemize}

Il moderatore conferma: come curatore dell'\textbf{area indie} di \textbf{Lucca
Comics \& Games}\footnote{Nella trascrizione ``Lukaku Vixen Games''. Segue un
riferimento poco chiaro (``Excalibur showcased some international indie games''),
non ricostruibile.} di quell'anno, ha ricevuto da alcuni studi internazionali demo
preparate \textbf{appositamente} per lo showcase: chi passa vuole giocare la parte
migliore del gioco in poco tempo.

\paragraph{Demo per eventi e demo online.} Un'altra domanda: meglio una demo corta
per gli eventi e una più lunga online? Qual è la durata ideale? Diego non farebbe
necessariamente due demo diverse: dipende dal \textbf{genere}. In generale
\textbf{30 minuti--un'ora} è una buona durata; per uno strategico profondo, online,
più vicino all'ora; per altri generi circa \textbf{30 minuti}, o anche meno se il
gioco lo permette. Il criterio è: quanto serve al giocatore per capire che cosa fa il
gioco e come sarà. Per il \textbf{playtest online} e la demo vera e propria si può
usare un modello diverso: a una fiera si può portare una beta o un prototipo che sia
piacevole da giocare, riservando il tempo di sviluppo alla demo online.

Il moderatore aggiunge che a una fiera, con altri studi intorno, bisogna
\textbf{catturare l'attenzione} di chi passa e guarda lo schermo per due o tre
secondi. Trenta minuti in presenza non sono troppi? Diego: è una stima
approssimativa; se il gioco si mostra in modo significativo in \textbf{15 minuti},
va bene. Non tutti i visitatori sono il pubblico di riferimento: chi è interessato
resterà di più, chi non lo è giocherà pochi secondi e se ne andrà, ``ed è meglio
così''. L'importante è che il \textbf{loop sia chiaro}, giocabile e ben eseguito.

\subsection{Appena laureati: un gioco grande o tanti piccoli?}

Con sei mesi o un anno liberi dopo la laurea, meglio un solo gioco o tanti piccoli?
Diego privilegia i \textbf{progetti piccoli}: conosce persone che hanno fatto una
serie di piccoli giochi, pubblicandoli, imparando, iterando e passando al successivo.
Quattro giochi gli sembrano troppi; con un anno, un paio di giochi è un buon
compromesso; con sei mesi, uno o al massimo due.

Il pericolo più grande è voler fare subito il \textbf{progetto dei sogni}: un gioco
che richiederà anni e molti soldi difficilmente riesce come primo progetto. Non
bisogna rinunciarvi, ma prima fare uno o due giochi più piccoli, comunque
appassionanti, per capire come si \textbf{costruisce, pubblica e promuove} un gioco.
Diego vede pagine Steam con la scritta ``coming soon'' da quattro o cinque anni: è
molto meglio un gioco \textbf{effettivamente pubblicato}, anche con poche recensioni
ma positive. ``Non farete \emph{Grand Theft Auto} o \emph{Call of Duty} come primo
gioco'', conclude, tra le battute del pubblico.

\subsection{Quando un evento va male}

Come reagire se un evento va male, con poca attenzione, poco feedback e pochi
incontri? Diego: può succedere che un evento deluda le aspettative; ciò che
\textbf{non deve succedere} è restare senza feedback e senza sapere perché.
\begin{description}
  \item[Fiere per il pubblico.] Alla fine della partita, passare uno o due minuti
        con chi ha giocato a chiedere com'è andata, e dare anche un piccolo
        \textbf{questionario}, perché la domanda diretta introduce un bias.
        Soprattutto \textbf{osservare} mentre giocano, invece di consegnare il
        gioco e chiacchierare con gli amici: dove le persone \textbf{si sporgono in
        avanti}, ``come diceva Jonas ieri''\footnote{Si veda la
        sezione~\ref{sec:c2-5-magia}.}, dove smettono di giocare, dove trovano il
        gioco troppo difficile. Confrontare ciò che si è osservato con ciò che le
        persone dicono: spesso dicono ``bellissimo'' ma hanno smesso perché
        frustrate. Il feedback è l'unica cosa davvero sotto il proprio controllo. Se
        l'evento va male, si capiscono le ragioni (dove si smette, poche wishlist),
        si \textbf{itera} sul gioco e si riprova a un altro evento dopo circa sei
        mesi.
  \item[Eventi B2B.] Un errore frequente: deck da 20 slide, incontro da 20
        minuti, 20 minuti di pitch senza lasciare spazio al publisher. Bisogna usare
        \textbf{al massimo metà del tempo} per presentarsi e presentare il gioco,
        lasciando il resto al \textbf{feedback}. Non tutti lo danno, ma molti sono
        felici di aiutare; anche nell'email di follow-up conviene chiedere un
        parere, per capire che cosa migliorare nel pitch, nella presentazione o nel
        gioco.
\end{description}

\subsection{Le candidature aperte}

In chiusura Diego ricorda che sono aperte le candidature per la \textbf{Official
Selection} (40 giochi) e per il mercato B2B, che abbina circa
\textbf{30}\footnote{La trascrizione riporta ``about 3 studios'', con ogni
probabilità un numero troncato.} studi a \textbf{40} publisher e investitori, sul
sito del festival, e si rende disponibile per domande dopo l'incontro.

% ---------------------------------------------------------------------
\section*{Punti chiave del capitolo}
\addcontentsline{toc}{section}{Punti chiave del capitolo}

\begin{itemize}
  \item Ogni evento ha un costo: valutare il \textbf{ritorno sull'investimento} e
        partire dagli eventi locali.
  \item Fiere per il pubblico: massima esposizione, meglio se con una
        \textbf{componente online} (showcase su Steam). Eventi B2B: imparare e
        incontrare persone, meglio se con una piattaforma di incontri.
  \item Per uno studio piccolo, meglio eventi europei di medie dimensioni che i
        giganti come Gamescom e GDC.
  \item Candidature e pitch: deck brevi (\textbf{max 10 slide}), immagini e
        gameplay, una slide di sintesi con genere, durata, uscita, budget, motore;
        se possibile, far \textbf{provare il gioco} durante l'incontro.
  \item Follow-up dopo una settimana, poi pazienza; lasciare metà dell'incontro al
        \textbf{feedback} del publisher.
  \item Due o tre eventi l'anno bastano: il costo opportunità è il tempo tolto al
        gioco.
  \item Partire da ciò che si ha da dire, non dalle tendenze; iniziare con progetti
        piccoli e \textbf{pubblicarli}.
  \item Demo: un loop completo, senza partire dal tutorial; osservare chi gioca e
        raccogliere feedback strutturato.
\end{itemize}
